\section{Real-valued inputs, positive witnesses, and examples}
\label{sec:extensions}

\subsection{Real-valued approximation functions with rational output covers}

The rational-valued restriction in the finite metric statement is convenient
for exact endpoint calculations.  It can be removed at the represented-input
level without adding an internal sort of real numbers.  Suppose a fixed
primitive recursive function $a(q,s)\in\Q$ names a nonnegative value
$\psi(q)$ with $|a(q,s)-\psi(q)|\le2^{-s}$.  Assume monotonicity of these
named values and a primitive recursive unit-tail modulus $D$ as before.
Every instance of these assumptions can be expressed using inequalities
between rational names.

At scale $Q=2^k$, set
\[
 e_k=2^{-(2k+4)},\qquad
 b_k=\max\{0,a(2^k,2k+4)-e_k\}.
\]
Then
\[
 0\le b_k\le\psi(2^k),\qquad
 0\le\psi(2^k)-b_k\le2e_k.
\]
Use $t_k=\min(2^k b_k,1)$ and the same rational radii
$r_k=t_k/(8\,4^k)$.  These intervals are contained in the original
approximation intervals because $b_k\le\psi(2^k)\le\psi(q)$ for
$q\le2^k$.  Monotonicity of $b_k$ itself is not required.

For a starting scale $K$, put $s=2^K$, $u=D(s)$ and $v=D(u+1)$.
The two disjoint blocks $s,\ldots,u$ and $u+1,\ldots,v$ each have
mass at least one, so the total mass through $v$ is at least two.
For $B=\lceil\log_2 v\rceil$, dyadic condensation yields
\[
 \sum_{k=K}^B2^k\psi(2^k)\ge2.
\]
The loss caused by the rational lower approximations is bounded by
\[
 \sum_{k=K}^B2^k\bigl(\psi(2^k)-b_k\bigr)
 \le\sum_{k=K}^B2^{-k-3}\le\frac14.
\]
Thus $\sum2^k b_k\ge7/4$, and capping at one still gives
$\sum t_k\ge1$.  Every subsequent counting, overlap and amplification
argument is unchanged.  The output is a rational list of intervals,
even when the original target radii are irrational.

For convergence, suppose the named values have a modulus $C(s)\ge1$
with all finite tail sums bounded by $2^{-s}$.  For a requested $m$ use
the rational upper envelope
\[
 u_m(q)=\max\{0,a(q,q+m+4)+2^{-(q+m+4)}\}.
\]
It satisfies $\psi(q)\le u_m(q)\le\psi(q)+2^{-q-m-3}$.
For $q\ge C(m+3)$, the open intervals of radius $u_m(q)/q$
centered at $p/q$, $0\le p\le q$, cover every original approximation.
Their unrestricted lengths, over any finite set of these denominators,
sum to at most
\[
 4\left(2^{-m-3}+\sum_{q\ge1}2^{-q-m-3}\right)
 =2^{-m}.
\]
Again the infinite geometric notation denotes a bound on every finite
prefix.  No positive lower bound for every $\psi(q)$ is needed in either
extension.

\subsection{An explicit extraction lemma for padded finite covers}

There is a constructive point in interpreting the null tests that deserves
a complete finite argument.  The negation of $a<x<b$ does not generally
give a decidable choice between $x\le a$ and $x\ge b$.  We avoid such
choices about real numbers altogether.

Let $U$ be a finite rational open interval list contained in $(0,1)$.
Compute its rational complement cells and the closed residual gaps by
sorting endpoints; include the singleton endpoints $0,1$ if necessary.
Pad every gap on both sides by a fixed positive rational $\delta$, and
write $V$ for that padded finite open cover.  Let $h>0$ be the least
positive difference between distinct endpoints of $U$ together with
$0,1$, and put
\[
 \varepsilon=\frac1{32}\min(\delta,h).
\]
The list contains $0,1$, so $h$ exists and is computed by a finite minimum.
Given a rational approximation $y$ to $x\in[0,1]$ with
$|x-y|<\varepsilon$, the following entirely rational case distinction
finds either an original interval containing $x$ with positive slack,
or a padded gap containing $x$ with positive slack.

First check whether $y$ is within $4\varepsilon$ of an endpoint $c$.
There is at most one distinct such endpoint because distinct endpoints
are at least $32\varepsilon$ apart.  If $c$ lies in an original open
interval, that interval's endpoints are at least $h$ away from $c$;
it therefore contains $x$, whose distance from $c$ is below
$5\varepsilon<h$.  If no original interval contains $c$, then $c$
belongs to a residual gap; its padded interval contains $x$ because
$|x-c|<5\varepsilon<\delta$.  These membership checks are comparisons
of rational endpoints.

If there is no such nearby endpoint, $y$ belongs to the interior of a
single endpoint cell and is more than $4\varepsilon$ from its boundary.
It cannot be outside $[0,1]$, since it is within $\varepsilon$ of $x$.
Then $x$ belongs to the same cell.  The cell is either covered by an
original interval, which supplies the desired interval, or is contained
in a residual gap, whose padding covers $x$.  Its coverage status is
read from the cell's rational midpoint.

Consequently, given the avoidance assumption that $x$ belongs to none of
the padded intervals, the second outcome is impossible and the finite
algorithm returns an original labelled interval containing $x$.  A
precision $s$ with $2^{-s}<\varepsilon$ is obtained by a bounded rational
calculation.  No search for a semidecidable real inequality, Markov's
principle, or excluded middle for real order has been used.

Apply this lemma to each residual list in the null-test construction of
Section~\ref{sec:finite-khintchine}.  Given a specific test index $m$ and
the assertion that $x$ avoids every listed interval of $V_m$, we obtain
a positive rational approximation beyond every prescribed denominator.
For fixed primitive recursive names, moduli and this index $m$, the
output witness is primitive recursive in the denominator threshold.
Iterating with the previous denominator as the next threshold produces
an increasing-denominator sequence of reduced approximants.

\paragraph{What is and is not asserted.}
The effective null test is a primitive recursive enumeration of rational
open intervals for each $m$, with a proved bound $2^{-m}$ on the sum of
the lengths of every finite sublist.  The positive conclusion takes an
explicit avoidance index as data.  We do not turn
$\neg(x\in\bigcap_m V_m)$ into $\exists m\,(x\notin V_m)$ by a
classical quantifier-negation rule.  In ordinary measure theory the
finite estimates imply the usual almost-everywhere theorem.  For arbitrary
classical monotone inputs, the same finite argument works with their
finite mass witnesses; that external interpretation does not assert
that those inputs have primitive recursive moduli.

\subsection{Concrete admissible inputs}

\paragraph{Golden ratio and the sharp coefficient.}
The continued-fraction digits $1,1,1,\ldots$ are primitive recursive.
Alternatively the name obtained by rational bisection of $z^2-5$ and
$I(Q)=4Q^2$ supplies the Hurwitz interface.  The optimality estimate has
a simple visible consequence: for $c=2/5$,
\[
 q^2<\frac{2c^2}{1-5c^2}=\frac85.
\]
Thus $|\phi-p/q|<(2/5)q^{-2}$ can hold only with $q=1$.
The exact coefficient $1/\sqrt5$ does give approximants beyond every
denominator, so retaining it in the constructive theorem is substantive.

\paragraph{Divergent harmonic examples.}
For any fixed positive rational $c$, take $\psi(q)=c/q$ and put
$r=\lceil2/c\rceil$.  The explicit elementary modulus
\[
 D(s)=2^r s-1
\]
works: on each of the $r$ blocks
$2^j s\le q<2^{j+1}s$, the sum of $1/q$ is at least $1/2$.
Hence the sum of $c/q$ from $s$ to $D(s)$ is at least $cr/2\ge1$.
In particular $c=1/4$ gives $D(s)=256s-1$ and target approximation
$|x-p/q|<1/(4q^2)$.  The finite coverage theorem applies to this
genuinely shrinking target, not only to large radii.

\paragraph{A convergent example.}
For $\psi(q)=q^{-2}$, the comparison
\[
 \frac1{q^2}\le\frac1{q(q-1)}=\frac1{q-1}-\frac1q\quad(q\ge2)
\]
gives a finite tail bound $1/(N-1)$.  Thus $C(m)=2^m+1$ is a valid
convergence modulus, and the target approximation is $|x-p/q|<q^{-3}$.
The null-cover conclusion follows with no monotonicity assumption beyond
what this particular example happens to satisfy.
